\section{来源审计：把增长故事改写成系统机制}

本讲由 Cursor CTO 与联合创始人 Sualeh Asif 主讲。公开视频在 2026 年 8 月 11 日核验时已经 private，仓库保留 1288 条时间戳字幕和官方封面。字幕记录了 2025 年 3 月的架构与事故经验；Cursor 2026 年的 secure indexing、privacy、CursorBench、Router 和 agent 文档只用于验证持久机制或说明后续演化，不倒写为课堂当时已经存在的能力。

\begin{warningbox}{规模数字的证据边界}
课堂提到过去一年约百倍增长、自研模型每日约一亿次调用、索引系统每日约十亿文档等数量。本文把它们保留为\textbf{讲者估计/课堂口径}，用于理解负载形状，不作为当前规模、商业排名或供应商关系的审计结论。
\end{warningbox}

\begin{importantbox}{全讲的三个闭环}
\begin{enumerate}
  \item \textbf{Context loop}：workspace change $\rightarrow$ incremental index $\rightarrow$ retrieval context；
  \item \textbf{Edit loop}：request $\rightarrow$ model/provider $\rightarrow$ plan/edit $\rightarrow$ fast apply $\rightarrow$ validation；
  \item \textbf{Reliability loop}：traffic/failure $\rightarrow$ telemetry $\rightarrow$ mitigation/migration $\rightarrow$ safer architecture。
\end{enumerate}
\end{importantbox}

\subsection{为什么 AI coding 不是一个模型 API}

AI coding 的用户体验由上下文 freshness、retrieval quality、模型能力、provider capacity、edit application 和 editor validation 共同决定。即使模型输出正确，索引过期、provider 排队、patch 应用慢或 diagnostics 卡顿，用户仍会感到系统失败。基础设施优化因此必须以完整 interaction loop 为对象。

\lecturefigure{01-coding-loop.png}{AI coding 产品是 workspace、context、model、apply 与 feedback 构成的闭环。}{本地字幕 00:00--05:20；概念重绘。}

\begin{knowledgebox}{读图：Latency budget 属于整条链}
把一次交互近似拆为
\[
T_{total}=T_{context}+T_{queue}+T_{model}+T_{apply}+T_{validate}.
\]
其中任一项的 p95/p99 都会破坏流畅感。Streaming 可以降低 time-to-first-token，却不能消除 context lookup、queueing、patch conflict 或 diagnostics 的尾延迟。
\end{knowledgebox}

\teachervoice{Asif 用“如果一亿次调用听起来不吓人，那它其实很吓人”提醒学生：增长改变的不只是容量，而是每个默认假设。曾经偶发的 retry、全量 rebuild 或人工修复，在高基数下都会成为持续负载。}

\subsection{本章小结}

Lecture 07 的主角不是 Cursor 品牌，而是一个高频 AI interaction system。后文每个架构选择都必须回答：它改善了哪一段 loop、增加了什么状态、故障时如何降级。

